% !TEX root = ../access.tex
%% ============================================================
%% SECTION 3 — METHOD
%% This file showcases the building blocks: equation, figure,
%% single-column table, and algorithm. Delete what you don't need.
%% ============================================================
\section{Proposed Method}
\label{sec:method}

Describe your method here. Refer to a numbered equation like
Eq.~\eqref{eq:example}, a figure like Fig.~\ref{fig:example}, a table like
Table~\ref{tab:example}, and an algorithm like Algorithm~\ref{alg:example}.

%% ---- Equation --------------------------------------------
\begin{equation}
y = f(\mathbf{x};\boldsymbol{\theta}) = \sum_{i=1}^{N} \theta_i\,x_i + b
\label{eq:example}
\end{equation}

%% ---- Single-column figure --------------------------------
% Put the image file in images/. Use width=\columnwidth for one column.
\begin{figure}[!t]
\centering
% \includegraphics[width=\columnwidth]{example_figure.png}
\caption{Single-column figure caption. Replace the commented
\texttt{\textbackslash includegraphics} line with your image.}
\label{fig:example}
\end{figure}

%% ---- Single-column table (stretched to column width) -----
% \extracolsep{\fill} pushes the table to fill the full column width.
\begin{table}[!t]
\caption{Single-Column Table Example}
\label{tab:example}
\setlength{\tabcolsep}{4pt}
\begin{tabular*}{\columnwidth}{@{\extracolsep{\fill}}lll@{}}
\hline
\textbf{Parameter} & \textbf{Value} & \textbf{Description} \\
\hline
Learning rate & $1\times10^{-3}$ & Optimizer step size \\
Batch size    & 64               & Samples per update \\
Epochs        & 100              & Training iterations \\
\hline
\end{tabular*}
\end{table}

%% ---- Algorithm -------------------------------------------
\begin{algorithm}
\caption{Example Procedure}
\label{alg:example}
\begin{algorithmic}[1]
\STATE \textbf{Input:} data $\mathbf{x}$, parameters $\boldsymbol{\theta}$
\FOR{each step $t = 1 \ldots T$}
  \STATE Compute output $y \leftarrow f(\mathbf{x};\boldsymbol{\theta})$
  \STATE Update $\boldsymbol{\theta}$ using the gradient
\ENDFOR
\STATE \textbf{Return:} $\boldsymbol{\theta}$
\end{algorithmic}
\end{algorithm}

% Tip: \FloatBarrier (from the placins package) forces all pending
% floats to be placed before this point — useful to keep a section's
% figures/tables from drifting into the next section.
